\begin{frame}{Motivação}
    \metroset{block=fill}

    \begin{block}{Problema}
        São dados um vetor $v_1, \dots, v_n$ de $n \leq 10^5$ inteiros e um inteiro $q \leq 10^5$. Teremos $q$ consultas, que podem ser de dois tipos:
            \begin{itemize}
                \item[(i)] recebe-se dois inteiros $l, r$ com $1 \leq l \leq r \leq n$, e deve-se imprimir a soma $v_l + v_{l+1} + \cdots + v_r$,
                \item[(ii)]recebe-se dois inteiros $i, x$ com $1 \leq i \leq n$ e $-10^9 \leq x \leq 10^9$, e deve-se atualizar $v_i = x$. Não é necessário imprimir valores nesse tipo de consulta.
            \end{itemize}
    \end{block}

    Vimos a técnica de soma de prefixos (psum). Qual é o problema dela?

    Ao atualizar o valor da posição $i$, é necessário atualizar os valores do vetor de psum da posição $i$ até a posição $n$. Em outras palavras, o valor $\operatorname{psum}[i]$ depende dos valores das posições $1, 2, \dots, i$.
    % \begin{center}
    %     \begin{tikzpicture}[every node/.style={circle, draw=black}]
    %         \node {1}[grow=0]
    %             child { node {2} 
    %                 child {node {3}
    %                     child {node {4}}}
    %             }
    %         ;
    %     \end{tikzpicture}
    % \end{center}
    
    % \begin{block}{Implementação da busca binária}
    %     \footnotesize
    %     \inputminted{cpp}{codigos/busca.binaria.cpp}    
    % \end{block}

\end{frame}
